\chapter*{Open Ends of the Theory}
\phantomsection\addcontentsline{toc}{chapter}{Open Ends of the Theory}
\markboth{Open Ends of the Theory}{}
\label{comp:open-ends}

The main manuscript lists its open problems in one table (\MT{app:claim-ledger}), each with a
settling condition: a class of models and two outcomes, a proof or a counterexample. The individual
chapters of this companion explain them where they arise. This chapter groups them by the question
a physicist would ask, says in which class each one is posed, and ends with the questions the program
cares about that the technical documents do not pose at all.

One feature of the settling conditions shapes everything below. Rule R16 of the authoring
specification requires both outcomes to be nontrivial within a stated class, so the problems are
posed in classes narrow enough to be decidable. Several of the central ones are posed on the
four-cycle kinetics of \MT{def:rg-formation-cycle-experiment} or on switching diffusions of that
type. Settling such a problem settles it for that class; it would ground the question at small size
but would not answer it for a network of hundreds of agents.

\section*{The parent's own dynamics}
\phantomsection\addcontentsline{toc}{section}{The parent's own dynamics}

Exact coarse-graining produces the law of a parent at every instant
(\MT{thm:rg-pointwise-parent-datum}). Whether that law evolves as an autonomous process, a parent
with dynamics of its own, is \MT{open:rg-pointwise-parent-dynamics}, which the introduction calls the
defining open problem of the phase. It asks whether a time-homogeneous Markov parent, or one carrying
a finite filter from a declared regular class, reproduces the native path law of the retained
coordinates within a total variation that vanishes as $\epsilon\to0$, both on windows of fixed length
and over the complete history. It is posed for switching diffusions of the four-cycle type. The
four-cycle continuation results of \cref{comp:ch14} answer special cases on the positive side: a
parent activated at a suitable time, after a coherence prefix, has vanishing discarded information and
a retained path law close to the native one on short windows. They do not settle the problem, because
they fix the activation event and the horizon. \MT{open:obs-selective-formation} asks the same
question for a block that communicates through learned low-rank interfaces: is there a parent of fixed
lower dimension whose path law of the communicated coordinates is asymptotically exact on the learning
time scale?

\section*{Formation that is not supplied}
\phantomsection\addcontentsline{toc}{section}{Formation that is not supplied}

The formation theorem of \cref{comp:ch13} relaxes a prepared state through a hierarchy of interaction
strengths supplied by hand. Four open problems ask for more. \MT{open:rg-formation-regime-sharpness}
asks whether the doubling recursion $\mathrm Q_\ell=2\mathrm Q_{\ell+1}+20$ of the strengths can be
relaxed to inner base $4/3$, a separation that is still doubly exponential in the depth, with constants
independent of the depth. \MT{open:rg-formation-adaptive-family} asks, on the
four-cycle, for parents with vanishing conditional and dynamical errors together with a margin,
bounded below, by which learning or the loop coordinate of the initial mean changes the law of the
first formed block; or for a proof that for almost every initial mean both differences vanish.
\MT{open:rg-formation-general} asks for repeated endogenous formation on general finite graphs with
temperature-independent data: a second coalescence that depends on the first by a margin, is held by
a restoring interaction, is observed on dynamical horizons, and has well-behaved parents; or a proof
that none exists in that class. \MT{open:full-graph-hierarchy} asks for the full package: a lineage of
at least two levels formed with probability bounded below from every initial law with a density and a finite
second moment, an intrinsic
partition selector, and a slow manifold on which the upper level feeds back consistently. These are
the questions on which the program's claim that collectives form through coupled inference rests.
None of the four takes a limit of many agents; the last three are posed at fixed finite size in the
low-temperature limit, while the regime-sharpness problem asks for constants uniform in the depth,
and hence in the number of agents of the family.

\section*{Selecting structure}
\phantomsection\addcontentsline{toc}{section}{Selecting structure}

No functional analyzed in the documents selects a nontrivial partition into meta-agents: the costs
studied either prefer an endpoint (every agent separate, or everything merged) or need an external
coefficient. \MT{open:restrict-intrinsic-criterion} and \MT{open:cg-natural-selector} ask for a
functional of the laws, natural under relabeling and compatible with the gauge action, not monotone
under refinement, with a nondegenerate optimizer, or for an impossibility theorem in a stated class.
\MT{open:rg-invariant-retained-subspace} asks whether deterministic block channels generically admit
retained coordinates beyond the trivial families of \cref{comp:ch11}. Structural dynamics of the graph
itself is open too: \MT{open:obs-edge-creation} asks for a normalized joint in which edges are created
and deleted by a threshold rule while the collective free energy does not increase, and
\MT{open:obs-multihead-attention} asks for an exact recursion for multi-head attention at the level of
source labels.

\section*{Convergence and time scales}
\phantomsection\addcontentsline{toc}{section}{Convergence and time scales}

The alternating scheme of inference and learning is not known to converge
(\MT{open:elbo-alternating-convergence}). The reduction of fast beliefs onto a slow manifold is proved
for held assignments; tracking along assignment paths with jumps is \MT{open:obs-slow-manifold},
posed on the four-cycle. Delayed feedback across learning episodes has an exact joint law, and whether
the per-episode accounting reproduces the aggregate update to first order is
\MT{open:full-graph-episode-accounting}. How a meta-agent's information clock compares with its
constituents' is \MT{open:hist-meta-clock-ratio}.

\section*{Foundations}
\phantomsection\addcontentsline{toc}{section}{Foundations}

Five open problems concern the measure-theoretic and geometric footing: likelihood families with no
common dominating measure (\MT{open:prob-nondominated-likelihood}), a normalized prior over law points
(\MT{open:prob-law-point-prior}), equivariant posterior versions for noncompact stabilizers
(\MT{open:cg-equivariant-versions}), global Fisher quotients (\MT{open:hist-retained-fisher-quotient}),
and correlated lifts beyond the finite parity witness (\MT{open:hist-parity-lift-scope}). The
equivariant-versions problem is open only when a law's stabilizer is noncompact and its posterior has
no weakly continuous version (\MT{prop:cg-equivariant-bayes-versions},
\MT{cor:cg-continuous-version-equivariant}). Noncompactness of $\GL(K)$ alone does not put the
program's Gaussian laws in the open case: under a faithful linear representation a nondegenerate
Gaussian law has a compact stabilizer, conjugate into $\mathrm O(K)$, so the open case concerns
degenerate laws, nonfaithful actions, and non-Gaussian fibers.

\section*{Dynamic links and their coarse laws}
\phantomsection\addcontentsline{toc}{section}{Dynamic links and their coarse laws}

The existence of a normalized generator prior is settled for every Lie group
(\MT{prop:link-generator-prior}); no compactness or finite Haar volume is needed. What remains
open is a finite-dimensional family of link laws that survives coarse multiplication
(\MT{open:link-coarse-closure}). For $\GL^+(K)$, $K\geq2$, the stated family must contain
the exponential images of all nondegenerate Gaussian generator laws and be closed under
convolution, with a weakly continuous injective parameterization by a connected open subset of a
finite-dimensional Euclidean space. The exact coarse law already exists; products can leave the
real exponential image, and local chart generators need not be Gaussian
(\MT{prop:link-coarse-generator-law}; \cref{comp:ch12}). This open item concerns useful
closure, not the existence of a parent pushforward.

Supplement~G poses two further questions (\cref{comp:gaussian-latent-links}). For unanchored
undirected Gaussian matter on connected multigraphs of cycle rank at least two, positive-definite
weights, zero information vector, and a product nondegenerate Gaussian generator base, is the
joint normalizer finite at a stated $K\geq2$ (\SG{open:glink-unanchored-normalization})?
The scalar case is finite, while an unanchored single cycle diverges. For the unregularized
expected comparison-record energy at $K\geq4$ with negative-trace rank-one optimum, can a chart
link having a nonsemisimple negative eigenvalue attain the chart infimum
(\SG{open:glink-chart-infimum})? Nonattainment is already proved for $K=2,3$ and for the
semisimple-negative class. Neither question is the earlier law-point-prior problem.

\section*{Contact with the world}
\phantomsection\addcontentsline{toc}{section}{Contact with the world}

Two items concern testability rather than mathematics. \MT{claim:physical-law-identification} asks
for held-out evidence, under a protocol fixed in advance, that an object of the coarse-graining
matches a named target system beyond a baseline. \MT{open:phil-testability} asks for a test at one
base point that could falsify nontrivial loop holonomy or the formation of a block, or a proof that
the records cannot distinguish them. The operational holonomy theorem of \cref{comp:ch10} shows that
records sensitive to network-loop holonomy exist in principle.

\section*{The supplements}
\phantomsection\addcontentsline{toc}{section}{The supplements}

Supplement~G keeps its own open rows: the scaling of a retained mass sector
(\SG{open:exp-mass-sector}), commutation of the information projection with coarse maps
(\SG{open:exp-projection-commutation}), monotonicity of Galerkin aggregation beyond one fixed precision
(\SG{open:grg-galerkin-monotonicity}), robustness of the anchored-star update
(\SG{open:obs-star-update-robustness}), Kron classification beyond finite encodings
(\SG{open:cg-kron-classification-beyond-finite}), a bridge to Bayesian renormalization
(\SG{open:grg-bayesian-bridge}), transfer of a model-space Fisher scale
(\SG{open:ig-model-space-fisher-transfer}), and stochastic inversion of aggregation
(\SG{open:grg-stochastic-inverse}), and the two link-sector questions described above.
\Cref{comp:sg1,comp:sg2,comp:gaussian-latent-links} explain them. Supplement~B holds the
open problems of the deferred program, among them continuum law theory
(\SB{open:prob-continuum-theory}), graph-to-base identification
(\SB{open:rg-cover-graph-base-identification}), compatible rescaling kernels on agent networks
(\SB{open:rg-rescaling-compatibility}), two-index limits and universality
(\SB{open:rg-two-index-limits}), and the typing of lattice sites (\SB{open:rg-lattice-typing});
\cref{comp:sb} explains them.

\section*{Questions the documents do not pose}
\phantomsection\addcontentsline{toc}{section}{Questions the documents do not pose}

Some questions the program cares about appear in none of the open-problem tables, either because
they were deferred by decision or because no one has yet posed them in a decidable form. Listing them
is not a criticism of the documents, which state their scope; it is a map of where the theory stops.

None of the main manuscript's open problems takes the limit of many agents. The limits taken are low
temperature, slow learning, and small separation parameters; two problems concern size without a
limit (uniformity in the depth for the regime-sharpness problem, and the dependence of a clock ratio
on the number of agents in \MT{open:hist-meta-clock-ratio}). The main manuscript's only large-size
statements are the retained-label Fisher theorems of \cref{comp:ch09}. They assume that the log
masses of a finite labeled mixture grow in proportion to an extensive size such as the number of
agents, and derive the Fisher geometry from that assumption; near a tie between labels the natural
coordinate is $N(\theta-\theta_0)$ (\MT{thm:rg-thermo-label-crossover}), the form of the finite-size
rounding of a first-order transition, here in a retained label. They do not derive the extensive
growth from a network of interacting agents, and they say nothing about formation. A thermodynamic
limit of the agent network at one base point, with nonabelian and noncompact links, is therefore
neither proved nor posed there; Supplement~B holds the infinite-network and asymptotic apparatus.
For the same reason, phase transitions, condensation, and critical slowing down are neither ruled in
nor ruled out as formation mechanisms: at finite size, with a unique law-level dynamics, they cannot
occur in the sense of statistical mechanics, and the main manuscript contains no statement about
formation in the limit in which they could.

A normalized law on noncompact links is supplied by the Gaussian-reference generator construction
of \cref{comp:lie-links}. The Haar-reference law remains outside the noncompact scope of
\MT{ch:geometry}; the chosen generator metric and prior are substantive model data. The
holonomy-conditioned parent score of \cref{comp:ch10} is infinite almost surely on nonforest
connected blocks under absolutely continuous $\mathfrak{gl}(K)$ generator laws
(\MT{cor:link-holonomy-parent-almost-sure}). That score is one particular construction, a projection of
marginal laws onto a holonomy-fixed family, not the exact coarse channel, whose parent law always
exists. Using this fixed nondegenerate Gaussian parent family on the stated $\GL^+(K)$ blocks
requires restricted generator support or a different parent family. Compact-unitary links
are outside that negative result.

Base-manifold transport, the setting of the program's original picture of agents spread over a base
with attention varying smoothly over overlaps, is deferred by the author's decision until the
single-point theory is settled; its apparatus and open problems are in Supplement~B.
